\chapter{Electrical \& Energy Storage Systems}
\label{ch:electrical_energy}

Dieses Kapitel behandelt die elektrotechnische Architektur, die Energiespeicherkomponenten sowie die ausfallsicheren Versorgungsnetzwerke des VTOL-Systems. Um die Betriebssicherheit des Gesamtsystems unter extremen atmosphärischen Bedingungen zu gewährleisten, wird eine strikte physische und galvanische Systemtrennung zwischen den Hochstrom-Antriebskreisen und den sensitiven Niedervolt-Avioniksystemen realisiert.

\noindent Der Fokus der Auslegung liegt auf der Implementierung eines intelligenten Batterie-Management-Systems (BMS) zur permanenten Einzelzellenüberwachung sowie auf der Definition konstruktiver thermischer Schutzmaßnahmen für den alpinen Worst-Case-Einsatz. Ergänzt wird die Energiearchitektur durch ein dual-redundantes Bus-Backup-System, welches bei einem totalen Ausfall der Hauptstromversorgung die unterbrechungsfreie Speisung aller flugsicherheitskritischen Steuer- und Aktuatorsysteme autonom sicherstellt.

\section{Power Distribution Architecture (High- vs. Low-Volt)}
\section{Battery Management System (BMS) \& Thermal Management}
\section{Emergency Power (Bus-Backup)}
